\section{Computational Implementation}

The numerical method was implemented in Python using a modular programming structure. Rather than organizing the solver as a single script, the implementation is divided into a collection of independent functions responsible for initialization, boundary-condition treatment, numerical solution, residual evaluation, grid generation, and post-processing. This modular design improves code readability, facilitates debugging, and provides a flexible framework for extending the solver to more advanced numerical methods in future projects.

Table~\ref{tab:functions} summarizes the principal functions comprising the solver.

\begin{table}[H]
\centering
\caption{Major functions of the numerical solver.}
\label{tab:functions}

\begin{tabularx}{\textwidth}{l X}
\toprule
\textbf{Function} & \textbf{Purpose} \\
\midrule

\texttt{initialize\_temperature()} 
& Constructs the initial temperature field and initializes ghost cells around the computational domain. \\

\texttt{apply\_boundary\_conditions()} 
& Updates ghost-cell values to enforce Dirichlet boundary conditions. \\

\texttt{residual\_norm()} 
& Computes the $L_2$ norm of the residual field for convergence monitoring. \\

\texttt{gauss\_seidel()} 
& Performs the iterative Point Gauss-Seidel solution of the discretized Laplace equation. \\

\texttt{create\_grid()} 
& Generates the structured computational grid for post-processing using \texttt{meshgrid}. \\

\texttt{plot\_contours()} 
& Compares numerical and analytical temperature distributions using contour plots. \\

\texttt{plot\_error\_field()} 
& Visualizes the spatial distribution of numerical error. \\

\texttt{plot\_line\_cuts()} 
& Extracts and plots horizontal and vertical temperature profiles for validation. \\

\texttt{plot\_residual\_history()} 
& Displays convergence history as a function of iteration count and computational time. \\

\bottomrule
\end{tabularx}

\end{table}

%%%%%%%%%%%%%%%%%%%%%%%%%%%%%%%%%%%%%%%%%%%%%%%%%%%%%%%%%%%%%

\subsection{Solution Algorithm}

The overall computational procedure consists of the following steps:

\begin{enumerate}
\item Specify the numerical parameters, including the mesh resolution and convergence tolerance.
\item Initialize the computational domain by constructing the temperature field together with the surrounding ghost cells.
\item Apply the prescribed Dirichlet boundary conditions through the ghost-cell values.
\item Repeatedly perform Gauss-Seidel iterations until convergence. Each iteration consists of
    \begin{enumerate}
        \item updating the interior control-volume temperatures,
        \item updating the ghost-cell values,
        \item computing the residual norm, and
        \item checking the convergence criterion.
    \end{enumerate}
\item Construct the computational grid for visualization.
\item Compare the numerical and analytical solutions using contour plots, error distributions, line-cut comparisons, and convergence histories.
\end{enumerate}

%%%%%%%%%%%%%%%%%%%%%%%%%%%%%%%%%%%%%%%%%%%%%%%%%%%%%%%%%%%%%

\subsection{Solver Implementation}

The Point Gauss-Seidel solver receives the initial temperature field, mesh resolution, and convergence tolerance as inputs. During each iteration, the temperatures of all interior control volumes are updated sequentially using Equation~\ref{eq:gs}. After completing one sweep of the computational domain, the ghost-cell values are updated according to Equation~\ref{eq:ghost}. The residual field is subsequently evaluated using Equation~\ref{eq:residual}, and its corresponding $L_2$ norm is computed from Equation~\ref{eq:l2}.

The residual norm, iteration number, and elapsed computation time are stored throughout the solution process to monitor convergence. The iterative procedure terminates once the prescribed convergence tolerance is satisfied, after which the converged temperature field is returned for post-processing.

%%%%%%%%%%%%%%%%%%%%%%%%%%%%%%%%%%%%%%%%%%%%%%%%%%%%%%%%%%%%%

\subsection{Grid Generation}

A structured computational mesh is generated by first constructing one-dimensional arrays containing the coordinates of the control-volume centers in each spatial direction. These arrays are subsequently converted into two-dimensional coordinate matrices using NumPy's \texttt{meshgrid()} function with the \texttt{indexing='ij'} option.

The \texttt{ij} indexing convention preserves the matrix indexing employed throughout the numerical solver such that

\[
X[i,j]=x_i,
\qquad
Y[i,j]=y_j.
\]

Consequently, the temperature field $T[i,j]$ corresponds directly to the physical location $(x_i,y_j)$ without requiring coordinate transposition during visualization. Maintaining this consistent indexing simplifies both the numerical implementation and the post-processing routines.

%%%%%%%%%%%%%%%%%%%%%%%%%%%%%%%%%%%%%%%%%%%%%%%%%%%%%%%%%%%%%

\subsection{Post-Processing}

Several post-processing techniques are employed to evaluate both the accuracy and convergence of the numerical solution.

Contour plots of the numerical and analytical temperature fields provide a qualitative assessment of the overall solution. An error contour illustrates the spatial distribution of the numerical error across the computational domain. Horizontal and vertical line cuts at selected locations enable direct quantitative comparison between the numerical and analytical solutions. Finally, the residual histories plotted as functions of both iteration number and computation time provide insight into the convergence characteristics and computational performance of the Point Gauss-Seidel solver.

